\section{既简单又复杂的系统：开源、数据墙与架构取舍}

上一章讲 K2 设计目标，本章看系统层取舍。为什么从闭源转向开源？为什么多模态不损伤“脑子”已经很好？为什么 Scaling Law 遇到数据墙后，FLOPs、RL、反馈和架构效率要重新平衡？这些问题说明，模型公司并不是只调一个旋钮，而是在复杂系统里不断取舍。

\begin{figure}[H]
\centering
\includegraphics[width=0.94\textwidth]{open-source-strategy.png}
\caption{从闭源到开源：模型与产品完成度、生态传播和商业路径共同影响开源选择。自制概念图，依据 00:54:08--01:05:00 对谈内容整理。}
\end{figure}

\begin{knowledgebox}{读图：开源是技术选择，也是生态选择}
闭源有利于控制产品体验和商业化，开源有利于生态扩散、开发者验证和信任建立。K2 开源意味着它既要证明基础模型能力，也要让 Agentic 生态围绕它生长。
\end{knowledgebox}

\subsection{Scaling Law、数据墙与 FLOPs}

上一节讨论开源和生态，本节回到模型训练的硬约束。杨植麟认为 Scaling Law 遇到数据墙是客观事实。高质量数据增长变慢以后，继续 scale 数据不再像过去那么简单；于是 FLOPs scaling、RL、环境 feedback、数据飞轮和架构效率的重要性上升。当前看，基于 FLOPs 的 scaling 仍然有效，但未来平衡点可能变化。

\begin{figure}[H]
\centering
\includegraphics[width=0.96\textwidth]{scaling-data-wall.png}
\caption{Scaling 与数据墙：数据墙之后，FLOPs、RL、反馈和架构效率重新平衡。自制概念图，依据 01:03:00--01:10:00 对谈内容整理。}
\end{figure}

\begin{warningbox}{数据墙不是训练停止}
数据墙不是说没有数据了，而是高质量、低噪声、可带来边际能力提升的数据变稀缺。模型公司仍然可以通过更好数据配比、合成数据、RL、工具环境和 test-time scaling 继续推进。
\end{warningbox}

\begin{knowledgebox}{数据墙后的四条路}
\begin{tabularx}{\textwidth}{p{0.20\textwidth}p{0.34\textwidth}X}
\toprule
路线 & 解决什么 & 代价 \\
\midrule
FLOPs scaling & 用更多计算继续压 loss & 成本高，依赖硬件供给。 \\
Data efficiency & 同样数据训练出更多能力 & 需要更好过滤、改写、配比和优化器。 \\
RL/feedback & 从环境反馈中创造新学习信号 & reward 噪声和安全问题更难。 \\
Architecture & 用新结构降低长上下文/推理成本 & 可能引入能力 bias。 \\
\bottomrule
\end{tabularx}
\end{knowledgebox}

\subsection{Linear Attention 与“智商”风险}

数据墙之后，架构效率会变得更诱人，本节讨论这种诱人的风险。访谈提到，很多 long context 架构会影响“智商”；纯粹 Linear Attention 可能因为架构 bias 影响模型能力。这个判断和 EP119 的 Attention 架构综述形成呼应：降低长上下文成本很重要，但如果架构改变牺牲了推理、记忆或表达能力，就可能影响模型上限。

\begin{figure}[H]
\centering
\includegraphics[width=0.94\textwidth]{linear-attention-iq-risk.png}
\caption{Long Context 架构与智商风险：Linear Attention 降成本，但 bias 可能影响模型能力。自制概念图，依据 01:17:17--01:20:00 对谈内容整理。}
\end{figure}

\begin{knowledgebox}{术语消化：架构取舍}
\begin{tabularx}{\textwidth}{p{0.22\textwidth}p{0.32\textwidth}X}
\toprule
路线 & 收益 & 风险 \\
\midrule
Full Attention & 表达力强、经验成熟 & 长上下文成本高，KV cache 压力大。 \\
Linear Attention & 成本低、长序列友好 & 架构 bias 可能影响能力。 \\
Sparse Attention & 降低计算，保留部分全局性 & 选择机制和硬件效率复杂。 \\
Hybrid & 效率和能力折中 & 比例、训练稳定和评测更复杂。 \\
\bottomrule
\end{tabularx}
\end{knowledgebox}

\subsection{基座模型公司与 Agent 产品公司的边界}

前面讨论模型内部的架构取舍，本节把问题推向产业边界。长期看，基座模型公司和 Agent 应用公司的边界并不固定。基座模型公司掌握模型能力、训练数据、API 和平台生态；Agent 应用公司掌握场景、工作流、用户入口和反馈数据。谁能把模型能力转成真实任务闭环，谁就能占据更高价值位置。

\begin{figure}[H]
\centering
\includegraphics[width=0.94\textwidth]{agent-product-boundary.png}
\caption{基座模型公司 vs Agent 应用公司：长期边界取决于模型、数据、工具和用户入口。自制概念图，依据 01:09:00--01:15:00 对谈内容整理。}
\end{figure}

\begin{importantbox}{边界判断}
如果 Agent 产品只是一层薄 workflow，边界会被基座模型公司吞掉；如果 Agent 产品拥有强场景、强数据回流和高频用户入口，它就可能反过来定义模型需求。
\end{importantbox}

\begin{knowledgebox}{商业模式对照：API、产品与生态}
\begin{tabularx}{\textwidth}{p{0.20\textwidth}p{0.34\textwidth}X}
\toprule
模式 & 优势 & 风险 \\
\midrule
API & 易规模化，开发者可接入 & token 价格竞争，离用户价值远。 \\
Agent 产品 & 更接近任务和用户入口 & 需要产品、工具、反馈和可靠性。 \\
开源生态 & 扩散快，信任和开发者参与强 & 商业捕获更难，支持成本高。 \\
闭源产品 & 体验可控，商业闭环清晰 & 生态扩散和外部验证受限。 \\
\bottomrule
\end{tabularx}
\end{knowledgebox}

商业模式还有一个成本侧问题：Agentic 任务通常比普通聊天消耗更多推理时计算。一个简单问答可能只需要一次模型调用；一个复杂 Agent 任务可能需要计划、搜索、工具调用、验证、重试和总结。这样会让“用户价值”和“推理成本”之间的差距变大：如果任务价值高，Agent 成本可以被吸收；如果任务只是低价值聊天，成本很容易压垮毛利。

\begin{knowledgebox}{Agent 成本结构}
\begin{tabularx}{\textwidth}{p{0.20\textwidth}p{0.34\textwidth}X}
\toprule
成本项 & 来自哪里 & 降本方式 \\
\midrule
模型调用 & 多轮规划、执行、反思 & 小模型路由、缓存、任务拆分。 \\
工具调用 & 搜索、代码、浏览器、API & 更好的工具协议和失败恢复。 \\
验证成本 & 测试、检查、人工确认 & 自动验证器和风险分级。 \\
上下文成本 & 长历史、文件、代码库、网页 & 记忆压缩和检索策略。 \\
\bottomrule
\end{tabularx}
\end{knowledgebox}

\subsection{本章小结}

K2 背后的系统取舍包括开源策略、数据墙、FLOPs scaling、架构效率和产品边界。它不是单点技术竞赛，而是模型公司在复杂系统中的多变量优化。

